% ------------------------------------------------------------------------------
% Chapter 3
% ------------------------------------------------------------------------------
\chapter{Implementation} 
\label{cha:implementation}

\section{FORMAL}
EBMC tool has been used to perform the bounded verification of this module.
Here the assumptions are the initial \texttt{rst\_n} low and the stability of inputs while \texttt{in\_valid} is high.
The assertions are as follows:
\begin{description}
    \item[\texttt{\_reset}]
        Verify that asserting a reset unconditionally forces the FSM back to the \texttt{IDLE} state on the next cycle with all registers cleared.
        \lstinputlisting[
            style=svlog, 
            linerange={155 - 163}, 
            firstnumber=155
        ]{gcd/gcd.sv}
    \item[\texttt{\_state}]
        Guarantees the FSM is mathematically restricted to the five valid encoded states, proving no illegal states can be reached.
        \lstinputlisting[
            style=svlog, 
            linerange={165 - 170}, 
            firstnumber=165
        ]{gcd/gcd.sv}
    \item[State Hold] (\texttt{idle\_idle}, \texttt{run\_run}, \texttt{done\_done}):
        Guarantee that the FSM retains its current state when no progression signal or abort is asserted.
        \lstinputlisting[
            style=svlog, 
            linerange={172 - 188}, 
            firstnumber=172
        ]{gcd/gcd.sv}
    \item[State Progression] (\texttt{idle\_done}, \texttt{idle\_run}, \texttt{run\_done}, \texttt{done\_idle}):
        Prove that the FSM correctly transitions to the next expected state when the appropriate progression flag is asserted.
        \lstinputlisting[
            style=svlog, 
            linerange={190 - 214}, 
            firstnumber=190
        ]{gcd/gcd.sv}
    \item[Output stability] (\texttt{outputs}):
        Verify that the results remain stable until the handshake happens.
        \lstinputlisting[
            style=svlog, 
            linerange={226 - 230}, 
            firstnumber=226
        ]{gcd/gcd.sv}
    \item[Handshakes] (\texttt{\_out\_valid}, \texttt{\_in\_ready}):
        Prove that the flags are never asserted spuriously and that they follow the FSM transitions.
        \lstinputlisting[
            style=svlog, 
            linerange={216 - 224}, 
            firstnumber=216
        ]{gcd/gcd.sv}
\end{description}

\section{HARDWARE}
The hardware checking happens inside the DUT interface and acts as a raw signal controller.
To do so, 5 properties with arguments to reuse them if necessary were created.
\lstinputlisting[
    style=svlog, 
    linerange={34 - 72}, 
    firstnumber=34
]{gcd/tb_uvm_gcd.sv}
These properties were used to make 7 assertions and 2 covers.
3 assertions are only used to control UVM drivers correctness.
\begin{description}
    \item[\texttt{chk\_reset}]
    Checks that in case of reset the DUT clear all of its outputs.
    \item[\texttt{chk\_early\_exit}]
    Checks that the DUT skips to the \texttt{RUN} state in case of early exit.
    \item[\texttt{chk\_out\_valid\_no\_drop}]
    Checks that \texttt{out\_valid} is not dropped unless a handshake happens.
    \item[\texttt{chk\_out\_data\_stable}]
    Checks that \texttt{gcd\_out} is stable while \texttt{out\_valid} is asserted.
\end{description}
The 2 covers simply make sure the handshakes, both input and output are covered.

\section{UVM}
The UVM architecture is organized as shown in the following diagram.
\begin{figure}[htbp]
  \centering
  \includesvg[width=\textwidth, inkscapelatex=false]{UVM.svg}
  \caption{UVM architecture of the DUT}
\end{figure}
